% ===========================================================================
%  离散数学与结构（I） 2026-10-15（第 3--4 节）课堂笔记
%  唯一分解与域扩张：从 UFD 到有限域
%  Compile:  xelatex main   (twice, for bookmarks and cross-references)
% ===========================================================================
\documentclass[cjk]{loom}

\newcommand{\N}{\mathbb{N}}
\newcommand{\Z}{\mathbb{Z}}
\newcommand{\Q}{\mathbb{Q}}
\newcommand{\R}{\mathbb{R}}
\newcommand{\C}{\mathbb{C}}
\newcommand{\F}{\mathbb{F}}
\newcommand{\Frac}{\mathrm{Frac}}           % 分式域
\newcommand{\chr}{\mathrm{char}}            % 特征

\runningthread{离散数学与结构 · 第四周（下）}

\begin{document}

\loomcover
  {唯一分解与域扩张：从 UFD 到有限域}
  {离散数学与结构（I）· 第四周（下）课堂笔记}
  {}
  {2026 年 10 月 15 日}

\begin{strand}
本讲脉络：在「整环」与「域」之间再分出三层——\keyword{UFD $\supset$ PID
$\supset$ ED}，一层比一层性质更好，而我们最关心的域上多项式环 $F[x]$
具有全部好性质。先复习整除与相伴，定义\keyword{素元}与\keyword{不可约元}，
给出\keyword{唯一分解整环}的定义；随即用 $\Z[\sqrt{-5}]$ 与 $\Z[\sqrt{10}]$
两个反例说明唯一分解并不自动成立。接着给出 UFD 的刻画
（\keyword{因子链条件} + \keyword{素性条件}），证明 PID $\Rightarrow$ UFD、
ED $\Rightarrow$ PID。应用上：PID 中 $(p)$ 极大当且仅当 $p$ 不可约，
由此用 $F[x]/(f)$ 构造域；再补充分式域的\keyword{泛性质}。后半讲转入
\keyword{域扩张}：扩域是基域上的线性空间，特征与\keyword{素域}，
有限域大小必为素数幂，扩张次数 $[E:F]$；$F[x]/(f)$ 中 $f$ 有根与
\keyword{嵌入定理}；最后引出\keyword{分裂域}。
\end{strand}

\clearpage
{\small\tableofcontents}

\clearpage
% ===========================================================================
\section{唯一分解整环（UFD）}\label{sec:ufd}

\subsection{概览：整环与域之间的三个层次}

环的谱系：
\[
  \text{环} \supset \text{交换环} \supset \text{整环} \supset \text{域}.
\]
本讲把「整环」到「域」之间再拆成三层，一层比一层性质更好：
\[
  \text{整环} \supset \text{唯一分解整环 (UFD)} \supset
  \text{主理想整环 (PID)} \supset \text{欧几里得整环 (ED)}.
\]
最重要的是：我们关心的环——\keyword{域上的多项式环 $F[x]$}——具有上述
所有好性质。

\subsection{单位、整除与相伴（复习）}

以下均在交换环 $R$ 中讨论。

\begin{definition}[单位（Unit）]
$a\in R$ 是\keyword{单位}，若存在 $a^{-1}\in R$ 使
$a\cdot a^{-1} = 1$。（例：$\Z$ 的单位只有 $\pm 1$。）
\end{definition}

\begin{definition}[整除]
$a\mid b$，若存在 $c\in R$ 使 $b = ac$。
\end{definition}

\begin{definition}[相伴（Associate）]
$a$ 与 $b$ \keyword{相伴}，若存在单位 $u\in R^\times$ 使 $a = bu$。
\end{definition}

若 $a$ 与 $b$ 相伴，显然 $a\mid b$ 且 $b\mid a$。反过来，在整环中也成立：

\begin{proposition}[互整除则相伴]
设 $R$ 为整环，$a,b\in R$。若 $a\mid b$ 且 $b\mid a$，则 $a$ 与 $b$ 相伴。
\end{proposition}

\begin{proof}
由条件存在 $c,d$ 使 $b = ac$，$a = bd$。于是 $a = adc$，即
$a\cdot 1 = a\cdot dc$。整环中有消去律，消去 $a$（设 $a\ne 0$）得
$dc = 1$，故 $c, d$ 都是单位。
\end{proof}

\subsection{素元与不可约元}

动机：整数中任何非零自然数都唯一地分解为素数之积（不计次序）。要推广到
一般整环，先定义「素」与「不可约」——分解只有限制在不可约因子上才谈得
上唯一。

\begin{definition}[素元（Prime Element）]
非零非单位元素 $p\in R$ 称为\keyword{素元}，若
\[
  p\mid ab \;\Longrightarrow\; p\mid a\ \text{或}\ p\mid b.
\]
\end{definition}

\begin{definition}[不可约元（Irreducible）]
非零非单位元素 $a\in R$ 称为\keyword{不可约元}（既约元），若
\[
  d\mid a \;\Longrightarrow\; d\ \text{与}\ a\ \text{相伴，或}\
  d\ \text{是单位},
\]
即 $a$ 只有平凡的因子。
\end{definition}

\begin{proposition}[素元不可约]
在整环中，素元一定是不可约元；反之不然（反例见第 \ref{sec:counter} 节）。
\end{proposition}

\subsection{唯一分解整环}

\begin{definition}[唯一分解整环（Unique Factorization Domain, UFD）]
整环 $R$ 称为\keyword{唯一分解整环}，若每个非零元 $r\in R$ 都可写成
\[
  r = u\, p_1 p_2 \cdots p_t,
\]
其中 $u\in R^\times$ 是单位，$p_i$ 是不可约元；且该分解在
\keyword{相伴与交换次序}的意义下唯一。
\end{definition}

\begin{remark}
\leavevmode
\begin{itemize}[leftmargin=1.3em,itemsep=1pt]
  \item 「唯一」必须在相伴意义下理解：$6 = 2\times 3 = (-2)\times(-3)$，
        但 $2$ 与 $-2$ 相伴，仍是同一个分解；
  \item 前面的单位 $u$ 用来涵盖 $r$ 本身是单位的情形（此时 $t = 0$）。
        若不愿意写 $u$，也可以把等号改为相伴关系。
\end{itemize}
\end{remark}

\begin{example}
$\Z$ 是 UFD（算术基本定理）。
\end{example}

% ===========================================================================
\section{反例：不是 UFD 的整环}\label{sec:counter}

\begin{definition}[{环的生成 $R[u]$}]
设 $R\subseteq R'$ 为环，$u\in R'$。$R[u]$ 表示 $R'$ 中包含 $R$ 与 $u$
的\keyword{最小子环}：
\[
  R[u] = \{\, r_0 + r_1 u + \cdots + r_t u^t : r_i\in R,\ t\ge 0\,\}.
\]
\end{definition}

于是 $\Z[\sqrt{-5}] = \{\,a+b\sqrt{-5} : a,b\in\Z\,\}$，
$\Z[\sqrt{10}] = \{\,a+b\sqrt{10} : a,b\in\Z\,\}$。它们都是整环：交换性
显然；无零因子因为它们是域（$\C$、$\R$）的子环。

\subsection{\texorpdfstring{$\Z[\sqrt{-5}]$}{Z[√-5]} 中分解不唯一}

\begin{example}[{$\Z[\sqrt{-5}]$ 中分解不唯一}]
\[
  6 = 2\times 3 = (1+\sqrt{-5})(1-\sqrt{-5}).
\]
各因子不可约的验证利用几何图像：把 $\Z[\sqrt{-5}]$ 看成复平面上的整点格。
模长满足 $|xy| = |x||y|$；离 $0$ 最近的非零格点是 $\pm 1$，故单位只有
$\pm 1$。找一个元素的因子只需在模长更小的格点中找——以原点为圆心、
$|a|$ 为半径的圆内只有有限个格点，逐一枚举即可验证 $2, 3,
1\pm\sqrt{-5}$ 均不可约。而 $2$ 与 $1\pm\sqrt{-5}$ 互不相伴（单位只有
$\pm 1$），两个分解之间没有相伴对应，故 $\Z[\sqrt{-5}]$ 不是 UFD。
\end{example}

这同时给出「不可约但不是素元」的例子：$2\mid 6 =
(1+\sqrt{-5})(1-\sqrt{-5})$，但 $2\nmid 1\pm\sqrt{-5}$。

\subsection{\texorpdfstring{$\Z[\sqrt{10}]$}{Z[√10]} 中分解不唯一}

\begin{example}[{$\Z[\sqrt{10}]$ 中分解不唯一}]
\[
  6 = 2\times 3 = (4+\sqrt{10})(4-\sqrt{10})\quad (16-10=6),
\]
\[
  15 = 3\times 5 = (5+\sqrt{10})(5-\sqrt{10})\quad (25-10=15).
\]
困难在于：$\{a+b\sqrt{10}\}$ 在实数轴上\keyword{稠密}，普通绝对值下
「模长更小的元素」有无穷多个，几何枚举法失效。解决办法是重新定义
「大小」——\keyword{范数（Norm）}：
\[
  N(a+b\sqrt{10}) = |a^2 - 10b^2|.
\]
范数是积性的，$N(x)N(y) = N(xy)$，可直接展开验证：
\[
  (a^2-10b^2)(c^2-10d^2) = (ac+10bd)^2 - 10\,(ad+bc)^2.
\]
因此 $x\mid y \Rightarrow N(x)\mid N(y)$。于是任何元素的潜在因子只有
范数严格更小的那些；范数取值为自然数，给定范数值的候选只有有限多种
（对 $a^2 - 10b^2$ 的给定值枚举 $a,b$）。由此同样可验证上面的因子均
不可约、两种分解互不相伴，故 $\Z[\sqrt{10}]$ 不是 UFD。
\end{example}

% ===========================================================================
\section{UFD 的刻画、PID 与 ED}\label{sec:pid-ed}

\subsection{UFD 的刻画：因子链条件 + 素性条件}

\begin{theorem}[UFD 的刻画]
整环 $R$ 是 UFD，当且仅当：
\begin{enumerate}[leftmargin=1.6em,itemsep=2pt]
  \item \keyword{因子链条件：}不存在无穷序列 $a_1, a_2, \dots$ 使每个
        $a_{i+1}$ 都是 $a_i$ 的\keyword{真因子}（$a_{i+1}\mid a_i$
        且二者不相伴）；
  \item \keyword{素性条件：}每个不可约元都是素元。
\end{enumerate}
\end{theorem}

\begin{proof}[UFD $\Rightarrow$ 两条件]
\leavevmode
\begin{itemize}[leftmargin=1.3em,itemsep=2pt]
  \item 若因子链条件不成立：$a_2\mid a_1$ 贡献一个不可约因子，
        $a_3\mid a_2$ 又贡献一个……$a_1$ 的分解将无穷长，矛盾；
  \item 素性：设不可约元 $p\mid ab$。把 $a$、$b$ 各自分解再拼起来，得到
        $ab$ 的一个分解；把 $ab/p$ 分解后与 $p$ 拼起来，得到 $ab$ 的
        另一个分解。由唯一性，$p$ 必与 $a$ 或 $b$ 的某个不可约因子相伴，
        故 $p\mid a$ 或 $p\mid b$。
\end{itemize}
\end{proof}

\begin{proof}[两条件 $\Rightarrow$ UFD]
与整数唯一分解的证明完全一样。因子链条件保证分解的\keyword{存在性}
（不断分解真因子的过程必须终止）。唯一性：设
\[
  p_1 p_2 \cdots p_r = q_1 q_2 \cdots q_s,\qquad
  p_i, q_j\ \text{均不可约}.
\]
由素性条件 $p_1$ 是素元，且 $p_1\mid q_1\cdots q_s$，故 $p_1$ 整除某个
$q_j$，不妨 $p_1\mid q_1$；$q_1$ 不可约而 $p_1$ 非单位，故 $p_1$ 与
$q_1$ 相伴。由消去律，$p_2\cdots p_r$ 与 $q_2\cdots q_s$ 相伴，递归做
下去即得两个分解在相伴与次序意义下相同。
\end{proof}

\subsection{主理想整环（PID）：PID \texorpdfstring{$\Rightarrow$}{=>} UFD}

\begin{definition}[主理想整环（Principal Ideal Domain, PID）]
由单个元素生成的理想 $(a) = aR$ 称为\keyword{主理想}。若整环 $R$ 中
\keyword{每个理想都是主理想}，则称 $R$ 为\keyword{主理想整环}。
\end{definition}

\begin{example}
$\Z$ 是 PID（上节课已证：其理想都是某个数的全体倍数）。
\end{example}

\begin{theorem}[PID $\Rightarrow$ UFD]
PID 是 UFD。
\end{theorem}

\begin{proof}
验证 PID 满足因子链条件与素性条件。

\keyword{因子链条件：}反设存在真因子链 $a_1, a_2, \dots$。由
$a_{i+1}\mid a_i$ 得 $(a_i)\subseteq(a_{i+1})$，即主理想升链
\[
  (a_1)\subseteq(a_2)\subseteq(a_3)\subseteq\cdots.
\]
其并 $I = \bigcup_i (a_i)$ 仍是理想（升链之并：对加减封闭显然；$x\in I$
属于某个 $(a_i)$，乘任何环元素后仍在 $(a_i)\subseteq I$ 中）。由 PID，
$I = (b)$ 是主理想。而 $b\in\bigcup_i(a_i)$，故存在 $i$ 使
$b\in(a_i)$，于是 $I = (b)\subseteq(a_i)$，即 $I = (a_i)$——链在第
$i$ 步后停止增长，$(a_i) = (a_{i+1}) = \cdots$，与「真因子链」矛盾。

\keyword{素性条件：}设 $p$ 不可约，$p\mid ab$。由 PID，$(a,p) = (s)$
是主理想。由 $p\in(s)$ 得 $s\mid p$；$p$ 不可约，只有两种可能：
\begin{itemize}[leftmargin=1.3em,itemsep=2pt]
  \item $s$ 与 $p$ 相伴：此时 $(a,p) = (p)$，故 $a\in(p)$，即 $p\mid a$；
  \item $s$ 与 $1$ 相伴：此时 $(a,p) = (1)$，存在 $\alpha,\beta$ 使
        $1 = \alpha a + \beta p$。
\end{itemize}
对 $b$ 同理。若 $p$ 既不整除 $a$ 也不整除 $b$，则
$1 = \alpha a + \beta p$ 与 $1 = \alpha' b + \beta' p$ 移项相乘得
\[
  (1-\beta p)(1-\beta' p) = \alpha\alpha'\, ab.
\]
右边被 $p$ 整除；左边展开为 $1 - (\beta + \beta' - \beta\beta' p)p$，
其中含 $p$ 的项被 $p$ 整除，于是 $p\mid 1$，$p$ 是单位——与不可约元
非单位矛盾。故 $p\mid a$ 或 $p\mid b$。
\end{proof}

\subsection{欧几里得整环（ED）：ED \texorpdfstring{$\Rightarrow$}{=>} PID}

如何证明一个环是 PID？给它一个更好的性质——带余除法。

\begin{definition}[欧几里得整环（Euclidean Domain, ED）]
整环 $R$ 称为\keyword{欧几里得整环}，若存在\keyword{范数}
$N\colon R\to\N$，满足：
\begin{enumerate}[leftmargin=1.6em,itemsep=2pt]
  \item $N(a) = 0 \iff a = 0$；
  \item \keyword{带余除法：}对任意 $b\in R$ 与 $a\ne 0$，存在 $q, r\in R$
        使
        \[
          b = aq + r, \qquad N(r) < N(a).
        \]
\end{enumerate}
\end{definition}

\begin{example}[最重要的例子]
域 $F$ 上的多项式环 $F[x]$ 是 ED，取
\[
  N(f) = \deg f + 1 \quad (f\ne 0).
\]
任给 $f\ne 0$ 与 $g$，多项式长除法给出 $g = fq + r$ 且
$\deg r < \deg f$，故 $N(r) < N(f)$。（$\Z$ 也是 ED，范数为绝对值。）
\end{example}

\begin{theorem}[ED $\Rightarrow$ PID]
欧几里得整环是主理想整环。
\end{theorem}

\begin{proof}
设 $I\trianglelefteq R$，$I\ne 0$。从 $I$ 的非零元中取范数最小者
\[
  a \in \arg\min_{x\in I,\ x\ne 0} N(x)
\]
（$N$ 取值于自然数，良序性保证最小值取得到）。断言 $I = (a)$：任取
$b\in I$，带余除法给出 $b = aq + r$，$N(r) < N(a)$。此时
$r = b - aq\in I$；但 $a$ 已是 $I$ 中非零元范数的最小者，唯一可能是
$N(r) = 0$ 即 $r = 0$。故 $a\mid b$，$I = (a)$。
\end{proof}

\begin{corollary}[{$F[x]$ 集所有好性质于一身}]
域上多项式环 $F[x]$ 是 ED，因而是 PID、UFD，具有之前所说的所有好性质。
\end{corollary}

% ===========================================================================
\section{极大理想与构造域}\label{sec:maximal}

\subsection{PID 中的极大理想}

回忆上节课的判别法（含幺交换环中）：
\[
  M\text{ 是极大理想} \iff R/M\text{ 是域}; \qquad
  P\text{ 是素理想} \iff R/P\text{ 是整环}.
\]
（域是整环，故极大理想必是素理想。）这给出构造域的方法：找一个整环，
取极大理想 $\mathfrak{m}$，则 $R/\mathfrak{m}$ 是域。在 PID 中判断极大
理想非常容易：

\begin{theorem}[PID 中极大理想的判别]
设 $R$ 为 PID，$a\in R$ 非零且非单位，则
\[
  a\text{ 不可约} \iff (a)\text{ 是极大理想}.
\]
\end{theorem}

\begin{proof}[双向都用逆否]
\leavevmode
\begin{itemize}[leftmargin=1.3em,itemsep=2pt]
  \item 若 $(a)$ 不极大，则存在 $J\supsetneq(a)$；$R$ 是 PID，$J = (b)$。
        由 $(a)\subsetneq(b)$ 知 $b\mid a$ 且不相伴，即 $b$ 是 $a$ 的
        真因子——$a$ 可约；
  \item 若 $a$ 可约，取真因子 $b$，则 $(b)\supsetneq(a)$，$(a)$ 不极大。
\end{itemize}
\end{proof}

\begin{corollary}[模掉不可约多项式得到域]
设 $F$ 为域，$f\in F[x]$ 不可约，则
\[
  F[x]/(f)\ \text{是一个域}.
\]
\end{corollary}

\begin{remark}[统一的观点]
之前见过的域也可以统一看成这种「模掉一个理想」的形式：例如
$\Q(\sqrt{2})\cong\Q[x]/(x^2-2)$——模掉理想相当于「给 $x$ 加一个
限制」：理想中的元素都被视为 $0$，即 $x^2 - 2$ 被看作 $0$，于是 $x$
就变成了 $\sqrt{2}$。
\end{remark}

\subsection{分式域的泛性质（复习与补充）}

\block{复习（分式域）}
设 $R$ 为整环，所有「分数」 $\frac{a}{b}$（$b\ne 0$）模去等价关系
$\frac{a}{b}\sim\frac{c}{d} \iff ad = bc$，所得集合构成域 $\Frac(R)$。

\begin{theorem}[分式域的泛性质 / 最小性]
$\Frac(R)$ 是把 $R$ 扩成域的\keyword{最小}办法：设
$\varphi\colon R\to R'$ 是单同态，且对任意 $x\ne 0$ 都有
$\varphi(x)\in(R')^\times$（$R'$ 是域时此条件自动成立），则存在
\keyword{唯一}同态 $\tilde\varphi\colon\Frac(R)\to R'$ 使
\[
  \tilde\varphi\circ\iota = \varphi, \qquad
  \iota\colon R\hookrightarrow\Frac(R),\quad x\mapsto\tfrac{x}{1}.
\]
\end{theorem}

\begin{proof}[构造]
$\tilde\varphi$ 的定义别无选择：
\[
  \tilde\varphi\!\left(\frac{a}{b}\right) = \varphi(a)\,\varphi(b)^{-1},
\]
这成立是因为 $\varphi(b)$ 在 $R'$ 中可逆。再验证它确实是同态即可。
\end{proof}

% ===========================================================================
\section{域扩张}\label{sec:extensions}

\subsection{域扩张与扩域的线性空间结构}

\begin{definition}[域扩张（Field Extension）]
设 $F\subseteq E$ 都是域，称 $F$ 是 $E$ 的\keyword{子域}，$E$ 是 $F$
的\keyword{扩域}（扩张）。
\end{definition}

\begin{example}
\leavevmode
\begin{itemize}[leftmargin=1.3em,itemsep=1pt]
  \item $E = F[x]/(p(x))$（$p$ 不可约）：$F$ 可看作 $E$ 的子域
        （常数多项式的等价类构成与 $F$ 同构的子域）；
  \item $E = \Frac(F[x])$（有理函数域）。
\end{itemize}
\end{example}

\begin{remark}[扩域的线性空间结构]
扩域 $E$ 天然地是基域 $F$ 上的\keyword{线性空间}：数乘就是域中的乘法，
逐条代入即知线性空间公理都满足。
\end{remark}

\subsection{特征与素域}

\block{复习（特征，Characteristic）}
含幺环 $R$ 的特征 $\chr R$ 是使
$\underbrace{1+\cdots+1}_{n\text{ 个}} = 0$ 首次成立的最小正整数 $n$；
若永不成立则为 $0$。
\begin{itemize}[leftmargin=1.3em,itemsep=2pt]
  \item $\chr R = n\ne 0$：由 $1$ 反复加减得到的元素构成子环，同构于
        $\Z/n\Z$；
  \item $\chr R = 0$：该子环同构于 $\Z$；
  \item $R$ 是\keyword{整环}且特征非零时，特征必为素数 $p$
        （整环无零因子）。
\end{itemize}

\block{域的情形（素域，Prime Field）}
设 $F$ 为域：
\begin{itemize}[leftmargin=1.3em,itemsep=2pt]
  \item $\chr F = p\ne 0$：$F$ 有子域同构于 $\F_p = \Z/p\Z$；
  \item $\chr F = 0$：先有子环同构于 $\Z$；域中非零元都是单位，由分式域
        的泛性质，$F$ 含有子域同构于 $\Frac(\Z) = \Q$。
\end{itemize}
这个最小的子域（$\F_p$ 或 $\Q$）称为 $F$ 的\keyword{素域}。

\subsection{有限域的大小必为素数幂}

\begin{theorem}[有限域的大小]
设 $E$ 为有限域，则 $|E| = p^k$ 为某个素数 $p$ 的幂。
\end{theorem}

\begin{proof}
有限域的特征必有限，故特征为某个素数 $p$；于是 $E$ 含有素域 $\F_p$；
而 $E$ 是 $\F_p$ 上的线性空间，故
\[
  |E| = p^k, \qquad k = \dim_{\F_p} E.
\]
\end{proof}

\begin{remark}[自然的问题]
是不是所有素数幂 $p^k$ 都能取到？构造方法：从 $\F_p$ 出发，取
$\F_p[x]$ 模掉一个 $k$ 次不可约多项式 $f$：
\[
  \F_p[x]/(f).
\]
由上一节的推论这是一个域，且有 $p^k$ 个元素（等价类由次数 $<k$ 的余式
代表）。所以问题归结为：\keyword{任意次数的不可约多项式是否存在}
（下一讲证明）。
\end{remark}

\subsection{扩张次数 \texorpdfstring{$[E:F]$}{[E:F]}}

\begin{definition}[扩张次数]
对域扩张 $F\subseteq E$，记
\[
  [E : F] := \dim_F E,
\]
即 $E$ 作为 $F$ 上线性空间的维数。注意与群论的记号不同：$[G:H]$ 表示
陪集的个数（指数）。
\end{definition}

% ===========================================================================
\section{嵌入定理与分裂域}\label{sec:splitting}

\subsection{\texorpdfstring{$F[x]/(f)$}{F[x]/(f)} 中
\texorpdfstring{$f$}{f} 有根；嵌入定理}

\begin{remark}[观察：不可约多项式在商环中获得根]
设 $f\in F[x]$ 不可约，$E := F[x]/(f)$。
\begin{itemize}[leftmargin=1.3em,itemsep=2pt]
  \item $f$ 在 $F$ 上不可约，故在 $F$ 中无根（若 $\alpha$ 是根，则
        $x - \alpha$ 是 $f$ 的因子）；
  \item 但在 $E$ 中 $f$ \keyword{有根}：根就是 $\bar x$（$x$ 在商环中的
        像）——$f(\bar x) = \overline{f(x)} = 0$，因为 $f(x)\in(f)$，
        而理想中的元素在商环中就是 $0$；
  \item 于是 $f$ 在 $E$ 中变得可约：在 $E[Y]$ 中（改用 $Y$ 表示变量），
        一次因子 $Y - \bar x$ 整除 $f(Y)$。
\end{itemize}
\end{remark}

\begin{theorem}[嵌入定理 / 构造的唯一性]
设 $f\in F[x]$ 不可约，$E = F[x]/(f)$。设 $E'$ 是 $F$ 的扩域且 $f$ 在
$E'$ 中有根 $u$（即 $f(u) = 0$），则存在\keyword{保 $F$ 的单同态}
\[
  \varphi\colon E = F[x]/(f)\hookrightarrow E'
\]
（保 $F$：$\varphi$ 限制在 $F$ 上是恒等）。
\end{theorem}

\begin{proof}
定义 $\varphi$ 为\keyword{赋值映射}：把等价类 $\alpha$ 映到
$\alpha(u)$，即 $\bar x\mapsto u$。
\begin{itemize}[leftmargin=1.3em,itemsep=2pt]
  \item \keyword{良定义：}若 $\alpha' = \alpha + q\cdot f$ 是同一等价类
        的另一代表，则 $\alpha'(u) = \alpha(u) + q(u)f(u) = \alpha(u)$
        （因 $f(u) = 0$）；
  \item \keyword{同态：}和在某点的取值等于取值的和，积同理；
  \item \keyword{单射：}域之间的非零同态必为单射（核是理想，而域只有
        平凡理想）。
\end{itemize}
\end{proof}

\begin{remark}[意义]
任何使 $f$ 有根的扩域 $E'$ 都同构地包含 $F[x]/(f)$——它是
\keyword{使 $f$ 有根的最小扩域}。注意嵌入一般不唯一：$f$ 在 $E'$ 中
每个根都给出一个嵌入（$\bar x$ 映到哪个根）。
\end{remark}

\subsection{分裂域}

\begin{definition}[分裂（Split）]
设 $f\in F[x]$。若 $f$ 在扩域 $E$ 上完全分解（在 $E[x]$ 中写成一次
多项式之积），称 $f$ 在 $E$ 中\keyword{分裂}。
\end{definition}

\begin{example}
$\R$ 上并非每个多项式都分裂；扩到 $\C$ 后所有多项式都分裂。使 $f$
分裂的扩域不唯一：$x^2 - 2$ 在 $\Q$ 上不可约，在 $\R$ 与 $\C$ 上都
分裂。
\end{example}

\begin{definition}[分裂域（Splitting Field）]
使 $f$ 分裂的\keyword{最小}扩域称为 $f$ 在 $F$ 上的\keyword{分裂域}。
\end{definition}

\block{目标定理（证明见下一讲）}
\begin{enumerate}[leftmargin=1.6em,itemsep=2pt]
  \item 对任何素数幂 $p^k$，存在大小为 $p^k$ 的有限域，且在同构意义下
        唯一；
  \item 对任何域 $F$ 与 $f\in F[x]$，分裂域存在且在同构意义下唯一。
\end{enumerate}

\begin{strand}
\keyword{分裂域存在性的直观构造：}把 $f$ 在 $F[x]$ 中分解为不可约因子；
一次因子不用管；取一个次数 $\ge 2$ 的不可约因子，按上一小节做一次商
（加入它的一个根），$f$ 就劈出一个一次因子；若尚未分解完全，继续往上
扩，有限步后得到分裂域。
\end{strand}

\begin{example}[一次取商一般得不到分裂域]
取 $\Q$ 上三次不可约多项式：加入一个根后只进来一个根，剩下一个二次
因子仍不可约，需要再做一次二次扩张。分裂域是六次扩张：先三次再二次，
$3\times 2 = 6$。
\end{example}

\begin{remark}[扩张顺序不影响结果]
从 $F$ 出发可以先加这个根、也可以先加那个根；对两条扩张路径分别用
嵌入定理构造单同态，拼起来即得两个扩域之间的同构。
\end{remark}


\end{document}
